\chapter{Background} \label{ch:background}
This chapter introduces all the background concepts required to understand the protocol and the context in which it works.
\section{Network architectures} \label{sec:network-architectures}
Network architecture defines how nodes communicate and interact over a network.
The most common architectural models are introduced in the following subsections.
\subsection{Distributed systems} \label{subsec:distributed-systems}
A distributed system is a collection of independent nodes that communicate and coordinate over a network to appear as a single coherent system~\cite{van2002distributed}.
This model enables horizontal scaling and eliminates single points of failure, making it suitable for high availability applications.
\subsection{Client-Server} \label{subsec:client-server}
The client-server architecture is one of the most widely adopted
architectural models~\cite{van2002distributed}.
It involves two types of nodes:
\begin{itemize}
    \item \textbf{Server}: the node responsible for delivering content
    and services to clients, maintained by the service provider.
    \item \textbf{Client}: the node that communicates with the server
    to request a specific service or resource.
\end{itemize}
These two node types communicate with each other following a
request-response model, where the client initiates the interaction
and the server responds accordingly.
\newline
Some protocols based on Client-Server architectures are:
\begin{itemize}
    \item \textbf{HTTP}: Client sends an HTTP request and receives an HTTP Response~\cite{rfc9110}
    \item \textbf{DNS}: Client (Resolver) sends a DNS query to DNS server to resolve a name~\cite{rfc1034}
    \item \textbf{DHCP}: Client broadcasts a DHCP request over all the local network and if a DHCP server is listening it will handle the IP assignation~\cite{rfc2131}
\end{itemize}

\subsection{Peer-to-peer} \label{subsec:peer-to-peer}
Peer-to-peer architectures are adopted in contexts where a single server is not able to guarantee a service, where the decentralization is a required property for the model or where the goal is to remove a single point of failure.~\cite{van2002distributed}

In a peer-to-peer context all nodes act as both client and server, simultaneously.
Each node is called \texttt{peer}, and it is able to request and provide resources to other peers.

In peer-to-peer protocols there is a challenge: peer discovery,
the mechanism by which a node locates other peers to communicate with.
A common approach is centralized trackers, where a server maintains a list
of active peers, but this reintroduces a single point of failure.
Distributed hash tables (DHT) solve this by distributing peer information
across the network itself, removing any central dependency~\cite{maymounkov2002kademlia}.

An example of a peer-to-peer protocol is BitTorrent~\cite{cohen2003incentives}, widely used for distributed file sharing.
The peer-to-peer model also serves as the foundation for distributed storage systems such as IPFS, discussed
in the following subsection.
\subsection{IPFS} \label{subsec:ipfs}
IPFS (InterPlanetary File System) is a peer-to-peer distributed file system that enables efficient and reliable data distribution among all nodes~\cite{benet2014ipfs}.
Unlike traditional location based addressing, each piece of content is identified by its cryptographic hash, known as a Content Identifier (CID).
The protocol works in the following way: every node hosts and shares a portion of the data, when a client requests a specific CID, it uses the DHT to retrieve the data from the nodes that store it.
In addition, IPFS is a versioned file system and it supports integration with decentralized applications, offering a storage layer for blockchains and Web3 ecosystems.
\section{Blockchain}\label{sec:blockchain}
A blockchain is a growing distributed ledger that stores data in an ordered sequence of blocks~\cite{bitcoin}.
Each block contains a set of transactions, the hash of the previous block, and a timestamp; this structure creates a chain.
This chained structure makes the blockchain transactions resistant to alteration, because, when the transactions are registered,
an adversary $\mathcal{A}$ would need to produce a block with the same hash, which is infeasible in probabilistic polynomial time~\cite{bitcoin}.
\subsection{Ethereum}\label{subsec:eth}
\subsection{Smart Contract}\label{subsec:sc}
\subsection{Gas}\label{subsec:gas}
\section{Crowdsensing}\label{sec:crowdsensing}
\section{Related work and state of the art} \label{sec:related-work}